\section{The primitive solution for independently signed slopes}
\label{cf:critical}

Put $r_\ell=1/p_\ell$ and solve
\begin{equation}\label{cf:r-equations}
 V_\ell(\boldsymbol\chi)+r_\ell U_\ell(\boldsymbol\chi)
                       =-2\pi\ii r_\ell .
\end{equation}
At $\mathbf r=0$ the solution is $\boldsymbol\chi=0$.
Lemma~\ref{cf:hessian} and \eqref{cf:UV} give
\begin{equation}\label{cf:Jacobian}
 D_{\boldsymbol\chi}V(0)=-2\pi\ii\,\diag(\tau_1,\ldots,\tau_H).
\end{equation}
It is invertible because $\Ima\tau_\ell>0$.
The holomorphic implicit-function theorem gives a unique solution
$\boldsymbol\chi_*(\mathbf r)$ on one fixed polydisk in all the
independent reciprocal slopes. Shrink that polydisk so that its
image lies in the framed neighborhood of
Section~\ref{cf:peripheral}. For real signed integer $p_\ell$
this requires only a lower bound on each $|p_\ell|$.
There is no constraint on their ratios.

When $p_\ell\ne0$, \eqref{cf:r-equations} is exactly
\begin{equation}\label{cf:slope-equations}
 U_\ell+p_\ell V_\ell=-2\pi\ii .
\end{equation}
By \eqref{cf:UV} these are precisely the critical equations of
\eqref{cf:action}. Thus local uniqueness applies to the stated
critical point, not only to a separate system of period equations.

There is also a direct classical check that the filled auxiliary
groups have a genuine loxodromic direction. Separate evenness of
$\Phi_\eps$ gives holomorphic functions $\widetilde\tau_\ell$
such that
\[
 V_\ell=\widetilde\tau_\ell(\boldsymbol\chi)\,U_\ell,\qquad
 \widetilde\tau_\ell(0)=\tau_\ell .
\]
In fact divide $\partial_{\chi_\ell}\Phi_\eps$ by $\pi\ii\chi_\ell$;
the apparent singularity is removable. Shrink the neighborhood
so $\Ima\widetilde\tau_\ell>0$ throughout it.
At the solution of \eqref{cf:r-equations},
\begin{equation}\label{cf:sign}
\begin{split}
 \chi_{\ell,*}
  &=\frac{r_\ell}
          {\widetilde\tau_\ell(\boldsymbol\chi_*)+r_\ell},\\
 \Rea V_\ell
  &=\frac{2\pi r_\ell^2\,
          \Ima\widetilde\tau_\ell(\boldsymbol\chi_*)}
           {|\widetilde\tau_\ell(\boldsymbol\chi_*)+r_\ell|^2}
       >0 .
\end{split}
\end{equation}
The second formula holds for either sign of the nonzero real
$r_\ell$. Thus the longitude multiplier has nonunit modulus.
The sign is derived from the classical potential and the original
positive cusp ratio, without a source noncancellation argument.

\section{From a peripheral equation to the original smooth filling}
\label{cf:metric}

We now apply a metric result, rather than identifying an
exponential equation with a filled metric.
The following is the precise input
\cite[Proposition~6.1]{Marked}.

\begin{proposition}[Near-complete primitive completion, imported]
\label{cf:completion-input}
Let $M$ be an oriented complete finite-volume hyperbolic
three-manifold with torus cusps and complete holonomy $\rho_0$.
Fix primitive bases $(\mu_i,\lambda_i)$.
There is one neighborhood of $\rho_0$ with its peripheral fixed
flags, chosen before any slopes, with the following property.
Let $\rho$ be an actual representation in that neighborhood and
$U_i,V_i$ its small similarity logarithms in bounded frames
sending the flags to infinity. Suppose both multipliers are one
at each cusp to remain unfilled. At each cusp to be filled suppose
$a_iU_i+b_iV_i=\pm2\pi\ii$ with $\gcd(a_i,b_i)=1$.
Then $\rho$ descends to the original topological filling and
is its smooth complete finite-volume hyperbolic holonomy.
\end{proposition}

This statement supplies completion; it does not supply the
representation or solve its equations. Its hypotheses hold here
for the following concrete reasons.
The complete manifold is the actual $Y^\eps_E$ of
\cite[Theorem~1.2]{Geometry}. Its primitive bases are the traced
curves \eqref{cf:peripheral-words}. The rational face maps and
original cut relations prove an actual framed representation
near that holonomy, in bounded charts. Section~\ref{cf:peripheral}
proves that every main peripheral group is a translation group.
The equations at the auxiliary cusps are
\eqref{cf:slope-equations}; their slopes are primitive because
$(1,p_\ell)$ is a primitive integer vector in the original basis.
These facts verify every premise of
Proposition~\ref{cf:completion-input}.

For clarity, we spell out the winding step in that proposition
and its relation to the present markings. On a compact truncated
core, a finite partition of unity in the future hyperboloid
model averages the small change of complete developing maps:
\begin{equation}\label{cf:core}
 \begin{split}
 Y_\rho(x)&=\sum_{j,\gamma}
       \eta_{j,\gamma}(x)\rho(\gamma)\rho_0(\gamma)^{-1}x,\\
 E_\rho(x)&=\frac{Y_\rho(x)}
       {\sqrt{-\langle Y_\rho(x),Y_\rho(x)\rangle}} .
 \end{split}
\end{equation}
Here the $\eta_{j,\gamma}$ are equivariant lifts of a finite
partition on evenly covered balls. A positive sum of future
timelike vectors is timelike. Reindexing proves equivariance,
and on a compact fundamental part $E_\rho$ is $C^1$-close to
the complete identity map. It is therefore a positive local
diffeomorphism. This argument assumes no discreteness of
the deformed representation.

For one complete cusp use coordinates
$(w,s)\in\CC\times[s_0,\infty)$ with translation lattice
$\Lambda$. In the moving flag frame its peripheral maps are
$A_\lambda(z)=a_\lambda z+b_\lambda$.
Choose a smooth compactly supported partition satisfying
$\sum_{\lambda\in\Lambda}\eta(w-\lambda)=1$, and put
\begin{equation}\label{cf:link}
\begin{split}
 d_\rho(w)&=\sum_{\lambda\in\Lambda}
            \eta(w-\lambda)A_\lambda(w-\lambda),\\
 F_\rho(w,s)&=\bigl(d_\rho(w),e^{s+\xi_\rho(w)}\bigr),
\end{split}
\end{equation}
where $\xi_\rho$ is the real-linear extension of
$\lambda\mapsto\log|a_\lambda|$.
Then $d_\rho(w+\lambda)=A_\lambda d_\rho(w)$.
Its real Jacobian is positive on a compact fundamental
parallelogram, by closeness to the identity, and equivariance
propagates this everywhere. The Jacobian of $F_\rho$ is
$e^{s+\xi_\rho(w)}\det_{\RR}Dd_\rho(w)>0$.
Interpolation with \eqref{cf:core} on a compact collar remains
$C^1$-close and positive.
This constructs an actual nonsingular hyperbolic metric on
the original core and its ideal-end-deleted cusps.

At a cusp to be filled, \eqref{cf:sign} in particular gives
a multiplier different from one. Commuting affine maps have
a common finite fixed point $q$, in addition to infinity.
For any generator with $a\ne1$ it is $q=b/(1-a)$.
Both complete generator translations are nonzero, while their
multipliers tend to one. Thus $|q|\to\infty$ uniformly near
completeness, even if one generator has multiplier one.
The developed fundamental parallelogram and its fixed generator
paths remain in a bounded set and have bounded lengths.
They therefore miss $q$ in one sufficiently small neighborhood.
Equivariance extends this exclusion to the entire link cover.

On that simply connected cover choose
$f=\Log(d_\rho-q)$. The metric $|df|^2$ is nonsingular and
descends to a compact torus; it is complete. Its lifted
Euclidean developing map covers $\CC$, by lifting finite line
segments. Since the plane is simply connected, $f$ is a
global diffeomorphism and its periods form a rank-two lattice.
For a fixed generator path of developed length at most $L$,
inside $|d_\rho|\le R$,
\begin{equation}\label{cf:small-period}
 |\Delta f|\le \frac{L}{|q|-R}\longrightarrow0.
\end{equation}
Its exponential is the actual affine multiplier. Consequently
the increment equals its prescribed small logarithm, with
no extra $2\pi\ii$ integer. Thus the two actual periods are
$U,V$ before the potentially large coefficient $p$ is used.
This is why the fixed neighborhood controls all independently
long slopes.

Translate $q$ to zero. Periodicity bounds
$e^{s_0+\xi_\rho(w)}/|d_\rho(w)-q|$, so the end contains
a full epigraph $y>A|d_\rho-q|$.
Set
\[
 r=\frac{y}{|d_\rho-q|},\qquad
 R=\operatorname{arsinh}(1/r),\qquad
 \theta=\Ima f,\qquad
 X=\Rea f+\tfrac12\log(1+r^2).
 \]
Direct substitution gives
\begin{equation}\label{cf:tube}
 ds^2=dR^2+\cosh^2R\,dX^2+\sinh^2R\,d\theta^2 .
\end{equation}
The epigraph contains a whole punctured tube $0<R<R_0$.
A period $L$ translates $(X,\theta)$ by $(\Rea L,\Ima L)$.
The original slope $\mu+p\lambda$ has period $-2\pi\ii$,
exactly one angular turn. Its complementary primitive period
$V$ has nonzero real part, either by rank two of the period
lattice or directly by \eqref{cf:sign}.
Adding the axis gives a smooth solid torus, with precisely
the original slope as meridian. Its holonomy is the identity:
it fixes both $q,\infty$ and has multiplier one.
This proves group descent, rather than merely an exponential
relation. Reversing the angular sign reverses its parametrization,
not the metric orientation.

\begin{figure}[ht]
\centering\begin{tikzpicture}[x=1cm,y=1cm,>=Latex,font=\small,
 period/.style={very thick,->,blue!60!black},
 boundary/.style={thin,gray!65}]
 \begin{scope}[shift={(0,0)}]
  \coordinate (o) at (0,1.8);
  \coordinate (m) at (0,-.4);
  \coordinate (v) at (3,2.5);
  \coordinate (mv) at (3,.3);
  \fill[blue!5] (o)--(v)--(mv)--(m)--cycle;
  \draw[boundary] (o)--(v)--(mv)--(m)--cycle;
  \draw[period] (o)--(m) node[midway,left] {$m=-2\pi\ii$};
  \draw[period] (o)--(v) node[midway,above] {$V$};
  \draw[->,gray!65] (m)--(mv);
  \draw[->,gray!65] (v)--(mv);
  \fill (o) circle (1.2pt);
  \node[above left] at (o) {$0$};
  \node[below,align=center] at (1.5,-.65)
       {Logarithmic link\\$U=m-pV$};
 \end{scope}
 \draw[->,thick] (3.7,1.05)--(4.8,1.05)
      node[midway,above] {fill};
 \begin{scope}[shift={(5.3,1.05)}]
  \draw[gray!55] (0,-.7)--(3.8,-.7);
  \draw[gray!55] (0,.7)--(3.8,.7);
  \draw[gray!55] (0,0) ellipse (.25 and .7);
  \draw[gray!55] (3.8,0) ellipse (.25 and .7);
  \draw[very thick,orange!70!black,->] (-.15,0)--(4.2,0)
      node[right] {$X$};
  \draw[blue!65!black,thick,->]
       (1.2,-.7) arc[start angle=-90,end angle=235,
                     x radius=.25,y radius=.7];
  \node[above,blue!65!black] at (1.2,.85) {$\theta\bmod 2\pi$};
  \node[below,align=center] at (1.9,-1.7)
      {Smooth tube about $R=0$\\
       $V:(X,\theta)\mapsto(X+\Re V,\theta+\Im V)$};
 \end{scope}
\end{tikzpicture}

\caption{The actual logarithmic link and its smooth filling.
The period $m=U+pV=-2\pi\ii$ is the primitive meridian of the
added solid torus; $V$ is the original disk-longitude period.
The old basis is $U=m-pV$. Equation \eqref{cf:tube} turns
$V$ into an axial translation and rotation, and $m$ into
one full angular turn.}
\label{cf:filling-figure}
\end{figure}

At an unfilled main cusp the multipliers are one. The same
compact Euclidean-link argument now gives a complete translation
plane and a full horoball quotient in \eqref{cf:link}.
The middle region is compact; the added tubes and the remaining
horoball ends have finite volume. The metric is complete.
Its universal developing map covers simply connected hyperbolic
space and hence is an isometry. The filled topology and all
its markings are the ones described in \eqref{cf:slope}.
This completes the metric part of Theorem~\ref{cf:main}.

\section{The absolute action and the filled fundamental chain}
\label{cf:volume-proof}

\subsection{The correct polarization}
Define $G=2\pi\ii\Phi_\eps$. The exact period formula
\eqref{cf:UV} gives
\begin{equation}\label{cf:primitive}
 \dd G=\frac12\sum_\ell V_\ell\,\dd U_\ell,\qquad
 \Ima G(0)=2\pi F_0=\Vol(Y^\eps_E).
\end{equation}
This is the $V\,dU$ polarization. A potential with differential
$-\frac12\sum U\,dV$ differs by the boundary term
$\frac12\sum UV$. Its regulator correction cannot be copied
unchanged into \eqref{cf:primitive}.

Let $\mathcal V(\rho_{\boldsymbol\chi})$ be the signed
ideal-tetrahedron volume of the retained original fundamental
chain from Section~\ref{cf:deformation}.
The local peripheral differential
\cite[equation~(7.3)]{Marked} is
\begin{equation}\label{cf:volume-differential}
 \dd\mathcal V
 =\frac12\sum_\ell
    \bigl(\Rea V_\ell\,\dd\Ima U_\ell
                 -\Rea U_\ell\,\dd\Ima V_\ell\bigr).
\end{equation}
Its hypotheses are present here: the oriented ideal chain is
positive and nondegenerate near the complete point; actual
face maps cancel the internal boundaries; edge relations
keep their logarithmic sums fixed; and the primitive
peripheral basis has the orientation
$\Ima(\lambda/\mu)>0$ at completeness.
All main logs vanish because their peripheral actions are
translations, so they contribute zero.
The factor $1/2$ in \eqref{cf:volume-differential} is for
similarity logarithms; replacing them by eigenvalue logarithms
without adjusting the factor would change the formula.

For reference, that local identity follows by differentiating
the Bloch--Wigner volume on each oriented tetrahedron and
cancelling paired normal-face flows. The remaining cusp
coefficient rows have intersection pairing $-2$ on the
primitive meridian/longitude basis; the inverse tangent
pairing is $-1/2$. In \cite[Section~7.1]{Marked} the signs
are computed with the positively ordered model
$(0,\infty,1,z)$ in standard upper-half-space orientation.
The original complete cusp orientation
\eqref{cf:cusp} is precisely that convention.
No complete character-variety uniqueness is an input to
this local calculation.

Taking the imaginary part of \eqref{cf:primitive} shows that
the differential of
$\Ima G-\frac12\sum\Rea U_\ell\Ima V_\ell$
is exactly \eqref{cf:volume-differential}.
At completeness its value is the actual metric volume,
by \eqref{cf:primitive}; hence the constant on this connected
neighborhood is fixed:
\begin{equation}\label{cf:regulator}
 \mathcal V(\rho_{\boldsymbol\chi})
   =2\pi\Rea\Phi_\eps(\boldsymbol\chi)
          -\frac12\sum_\ell\Rea U_\ell\,\Ima V_\ell .
\end{equation}
This is an absolute equality on the original marked chain.

\subsection{The critical boundary term}
At a real integer filling tuple, write $\chi_\ell=u_\ell+\ii v_\ell$.
The critical equation gives
$V_\ell=2\pi\ii(\chi_\ell-1)/p_\ell$.
Direct substitution yields
\begin{equation}\label{cf:boundary-cancellation}
\begin{split}
 -\tfrac12\Rea U_\ell\,\Ima V_\ell
  &=-\frac{2\pi^2}{p_\ell}v_\ell(u_\ell-1)\\
  &=2\pi\Rea\left(\frac{\pi\ii}{2p_\ell}
                              (\chi_\ell-1)^2\right).
\end{split}
\end{equation}
Thus \eqref{cf:regulator} equals
$2\pi\Rea S_{\eps,\mathbf p}(\boldsymbol\chi_*)$.
It remains to identify the same ideal-chain volume with
the complete filled metric, rather than infer that conclusion
from a differential identity.

\subsection{Keeping the same boundary triangulation}
The comparison in \cite[Section~7.3]{Marked} applies to the
primitive completion proved above. We describe its chain
content to fix the absolute interpretation.
At a filled cusp conjugate the two fixed points to $0,\infty$.
On each triangle of the retained boundary triangulation
the conical ideal development has form
$(a(w),Yb(w))$, with $b>0$.
The functions and their derivatives are bounded on a finite
fundamental set. An equivariant axis map is
$(0,Y e^{\eta(w)})$, where $\eta$ has the real peripheral
logarithmic character. On that same triangulation use
\begin{equation}\label{cf:homotopy}
 H_s(w,Y)=\left(sa(w),
     Y\exp\bigl((1-s)\eta(w)+s\log b(w)\bigr)\right),
       \qquad 0\le s\le1.
\end{equation}
It is face-compatible and equivariant: horizontal coordinates
transform by the complex multiplier and both heights by its
modulus. The height is comparable to $Y$, horizontal
derivatives are bounded, and height derivatives are $O(Y)$.
The hyperbolic volume form
$y^{-3}\,dx_1\wedge dx_2\wedge dy$ bounds this boundary
homotopy's volume by $O(Y^{-2})$.

After the primitive slope has been killed, extend the same
axis-valued triangulated boundary through the solid torus.
The remaining cyclic peripheral group acts along the axis.
All added vertices lie on that axis, so every added
tetrahedron has zero signed volume.
Unfilled main cusp comparisons have the same vanishing
boundary estimate.
Stokes on finite truncations, followed by $Y\to\infty$,
therefore identifies the original chain with the filled
equivariant fundamental class.

Compare it with the developing map of the complete metric
from Section~\ref{cf:metric}. They represent this same class
and have the same remaining cusp flags. A common subdivision,
straightening and the same vanishing boundary estimate give
\[
 \mathcal V(\rho_{\boldsymbol\chi_*})
       =\Vol\bigl(Y^\eps_E(\mathbf p)\bigr).
 \]
Using an unrelated new triangulation with axis vertices
would not prove this equality; it is the retained boundary
homotopy \eqref{cf:homotopy} that supplies the comparison.
The usual metric volume is unambiguous also by the
volume-independence statement \cite[Lemma~7.1]{Marked}.
Together with \eqref{cf:boundary-cancellation} this proves
\eqref{cf:volume} and Theorem~\ref{cf:main}.

\section{The analytic interface and finite checks}
\label{cf:interface}

The classical potential has been defined from the cell functions
\eqref{cf:m}, \eqref{cf:g} and the exact meridian shear.
An application to a finite source must separately prove that
its selected limiting leaf action is this $f_\eps$, including
the parameter-dependent constant. Matching only its $x$
derivative or its complete Hessian is insufficient.
For example the exact classical exponential relation
$R^{-1}-1=e^{-\partial_x f_\eps}$ fixes a face shape,
but does not by itself identify a finite coefficient's
normalization. One must also compare the endpoint-normalized
leaf primitive and then eliminate the same internal wires.
Control of complementary terms and a nonzero whole filling
asymptotic are further analytic assertions. None is a
premise or conclusion of Theorem~\ref{cf:main}.

The construction-specific imported inputs are the complete
stationary point of \cite[Proposition~5.3]{Newton} and the
complete metric, actual cut and primitive bases of
\cite[Theorem~1.2 and Sections~4--7]{Geometry}.
The full complex continuation uses the latter paper's
rational face identities, not a claim that a nonplanar
deformed half ball is convex. The metric completion and
peripheral-volume facts imported from \cite{Marked} have
their actual hypotheses verified above. Standard holomorphic
implicit-function theory is used only at the two explicitly
invertible Jacobians. The usual Dehn-filling and
peripheral-volume framework provides historical context
\cite{NeumannZagier}; the original words and the absolute
normalization are retained in the calculation here.

The deterministic accompanying program \texttt{check.py}
verifies exact rational flag and multiplier identities,
the local Hessian coefficient, the real signed formula
\eqref{cf:sign}, and the boundary cancellation
\eqref{cf:boundary-cancellation}. Its negative controls
detect the missing factor two between eigenvalue and
similarity logs and an incorrect shear sign.
The record includes hashes of the actual source inputs.
These checks establish finite algebraic equalities only.
They do not prove the implicit-function theorem, the original
topological identification, metric completeness, or the
fundamental-chain comparison.

\begin{thebibliography}{9}
\bibitem{Geometry}
\emph{Pyramids, prisms, and complete geometry of Conway closures},
companion paper, Theorem~1.2 and Sections~4--7.
\href{https://sashakolpakov.github.io/murakamis_dream/papers/conway-geometry.pdf}
{Companion PDF}.

\bibitem{Newton}
\emph{Newton coefficients and positive sources for Conway tangles},
companion paper, Proposition~5.3.
\href{https://sashakolpakov.github.io/murakamis_dream/papers/newton-tangles.pdf}
{Companion PDF}.

\bibitem{Marked}
\emph{Circle patterns and primitive fillings of augmented links},
companion paper, Proposition~6.1, equation~(7.3) and Section~7.3.
\href{https://sashakolpakov.github.io/murakamis_dream/papers/marked-fillings.pdf}
{Companion PDF}.

\bibitem{NeumannZagier}
W.~D.~Neumann and D.~Zagier, Volumes of hyperbolic three-manifolds,
\emph{Topology} \textbf{24} (1985), 307--332.
\href{https://doi.org/10.1016/0040-9383(85)90004-7}
{doi:10.1016/0040-9383(85)90004-7}.
\end{thebibliography}
